\label{chap:introduction}

\section{Background and Motivation}
Autoregressive language models like GPT-2 use the Transformer architecture \cite{vaswani2017attention}, relying on scaled dot-product attention to connect tokens across a sequence:
\begin{equation}
\label{eq:standard_attention}
\text{Attention}(Q, K, V) = \text{softmax}\left(\frac{Q K^T}{\sqrt{d_k}}\right) V, \quad \text{where } \text{softmax}(z_i) = \frac{\exp(z_i)}{\sum_j \exp(z_j)}
\end{equation}

While dot-product attention works remarkably well, its normalization function, softmax, has a clear mathematical property: for any finite real input $z$, $\exp(z)$ is strictly positive. As a consequence, softmax produces a non-zero probability for every token in the sequence. The model must assign some attention mass to every word in its context window. Even when tokens are completely unrelated—such as punctuation marks, filler words, or noisy passages—they still receive a small fraction of the attention. Over several layers, this background leakage compounds, diluting relevant context representations and increasing the chance of factual hallucinations \cite{ji2023survey}. Furthermore, computing pairwise attention over all tokens forces an $O(S^2)$ memory and computation cost, even when most token pairs share no meaningful relationship.

\section{Core Idea: Multi-Attribute Continuous Logic}
In this work, we take a different perspective on token routing. Instead of treating tokens as unbounded vectors in Euclidean space $\mathbb{R}^D$ and measuring angular similarity, we treat tokens as collections of continuous attributes bounded in the unit hypercube $[0, 1]^D$. Under this representation, token interaction becomes a question of multi-attribute satisfaction and logical containment.

In classical logic, when two statements do not overlap, their conjunction evaluates to false. To bring this hard cutoff into a continuous neural network, we use the \textbf{nilpotent \L{}ukasiewicz triangular norm ($t$-norm)} \cite{klement2000triangular}, defined as $T_L(a, b) = \max(0, a + b - 1)$. If the sum of two attribute values $(a + b)$ is at or below $1$, the output evaluates to exactly zero ($0.0000$). With an activation threshold $\tau$, this operator serves as a differentiable gate: whenever semantic overlap between tokens is insufficient, the routing weight drops straight to zero.

\section{Problem Statement and Objectives}
This thesis asks a practical research question: \textit{Can continuous soft set representations in $[0, 1]^D$ combined with nilpotent triangular norms build a sequence router that achieves true exact-zero sparsity and shields against distractors without breaking language generation?}

We set four specific goals:
\begin{enumerate}[itemsep=0.5pt, topsep=1.0pt, parsep=0pt]
    \item Formulate a continuous lattice router using bounded embeddings, nilpotent \L{}ukasiewicz $t$-norms, and Jaccard containment.
    \item Train a 13.67M-parameter model matching a standard Transformer baseline on WikiText-2 and TinyStories under identical compute budgets.
    \item Evaluate distractor resistance using synthetic Needle-in-a-Haystack tests and test logical reasoning on multi-step Modus Ponens deduction chains.
    \item Analyze empirical trade-offs between exact-zero sparsity, generation speed, and next-token perplexity.
\end{enumerate}

\section{Summary of Key Findings and Outline}
Our experimental results highlight four main takeaways:
First, regarding perplexity, Soft L-Topology does not surpass standard Transformers (3.7455 loss / 42.33 PPL vs. 3.4866 loss / 32.67 PPL on TinyStories). Softmax produces smooth gradients that naturally benefit cross-entropy optimization.
Second, the model delivers on sparsity, sustaining an average of \textbf{55.06\% exact-zero routing} ($R_{ij} \equiv 0.0000$) throughout autoregressive generation.
Third, in Needle-in-a-Haystack evaluations with 16 injected distractors, it prunes 66.9\% of distractor connections to exact zero, reducing attention leakage by \textbf{49.4\% in absolute terms}.
Finally, on structured reasoning, it doubles Modus Ponens accuracy (40.0\% vs. 20.0\%) and achieves lower Masked Language Modeling loss (6.2041 vs. 6.2422) under heavy distractor noise.

The rest of the thesis is organized as follows: Chapter 2 discusses theoretical foundations; Chapter 3 describes the architecture and lattice KV-caching; Chapter 4 presents setup and baseline ablations on WikiText-2; Chapter 5 details evaluations on TinyStories and critical discussion; Chapter 6 summarizes conclusions and future directions; and Appendix A provides mathematical proofs and test specifications.
